\documentclass[11pt]{article}
\usepackage[T1]{fontenc}
\usepackage[utf8]{inputenc}
\usepackage{lmodern}
\usepackage[a4paper,margin=27mm]{geometry}
\usepackage{amsmath,amssymb,amsthm,mathtools,mathrsfs}
\usepackage{microtype}
\usepackage{booktabs,array,longtable}
\usepackage{enumitem}
\usepackage{xurl}
\usepackage{xcolor}
\usepackage[colorlinks=true,linkcolor=blue!45!black,citecolor=blue!45!black,urlcolor=blue!45!black]{hyperref}
\usepackage[nameinlink,noabbrev]{cleveref}
\allowdisplaybreaks
\emergencystretch=2em
\setlength{\parindent}{0pt}
\setlength{\parskip}{0.48em}

\newtheorem{theorem}{Theorem}[section]
\newtheorem{proposition}[theorem]{Proposition}
\newtheorem{lemma}[theorem]{Lemma}
\newtheorem{corollary}[theorem]{Corollary}
\theoremstyle{definition}
\newtheorem{definition}[theorem]{Definition}
\newtheorem{example}[theorem]{Example}
\theoremstyle{remark}
\newtheorem{remark}[theorem]{Remark}

\newcommand{\Z}{\mathbb Z}
\newcommand{\Q}{\mathbb Q}
\newcommand{\R}{\mathbb R}
\newcommand{\C}{\mathbb C}
\newcommand{\Pone}{\mathbb P^1}
\newcommand{\ES}{\operatorname{ES}}
\newcommand{\SL}{\operatorname{SL}}
\newcommand{\GL}{\operatorname{GL}}
\newcommand{\Sym}{\operatorname{Sym}}
\newcommand{\diag}{\operatorname{diag}}
\newcommand{\Pf}{\operatorname{Pf}}
\newcommand{\SNF}{\operatorname{SNF}}
\newcommand{\Li}{\operatorname{Li}}
\newcommand{\Spec}{\operatorname{Spec}}
\newcommand{\dd}{\,\mathrm d}
\newcommand{\eps}{\varepsilon}
\newcommand{\1}{\mathbf 1}
\newcommand{\mat}[4]{\begin{pmatrix}#1&#2\\#3&#4\end{pmatrix}}

\title{Divisor Channels, Prime Saturation, and Bézout Openings\\
\large An exact Erdős--Straus--Koide--period-discriminant extension}
\author{Technical research continuation from the cited canonical workbenches}
\date{4 September 2026}
\hypersetup{
 pdftitle={Divisor Channels, Prime Saturation, and Bezout Openings},
 pdfauthor={Technical research continuation from the cited canonical workbenches},
 pdfsubject={Exact Erdos-Straus, Koide, and period-discriminant mathematics}
}

\begin{document}
\maketitle

\begin{abstract}
Starting from the exact marked Erdős--Straus selector already identified with
an integral terminal-matrix slice of the proposed complex-$S^6$ construction,
we extend the arithmetic chain in four directions.  First, for every coprime
odd fixed-denominator shell, not only for prime denominators, we classify all
oriented reciprocal solutions by divisor channels $\delta\mid n^2$ and prove
a fixed-point-free complementary-channel involution.  Each channel has an
explicit terminal relation matrix of determinant $-\delta$, and its channel
count has an exact finite Fourier expansion over
$(\Z/R\Z)^\times$.  Second, primitive factor coordinates turn every channel
into the projective incidence $nr+\delta t\equiv0\pmod R$ and into an integral
Lorentz similitude of factor $n^2$.  Third, every primitive positive Koide
null ray supports infinitely many exterior prime certificates in every
admissible prime progression; the single diagonal ray $[1:1]$ alone captures
relative density one of the primes in every such progression.  Fourth, a
chosen Bézout companion opens each rank-one selector ray to an explicit
nondegenerate determinant-$(-\delta)$ coefficient matrix.  This matrix gives
a regular ambient period torus, is related to the terminal relation matrix by
a fully displayed $\GL_2(\Z)$ left--right morphism, and maps by an exact
orthogonal--homothetic map to the defect-one Koide hyperboloid.  That exact map
is proved to have no rational image point, so the existing integral defect-one
counting theorem cannot be silently reinterpreted as an Erdős--Straus count.
On the diagonal primitive ray, the number of exterior certificates for each
individual prime is computed exactly from the factorization of $(p+1)/2$.
\end{abstract}

\paragraph{Source status.}
The imported prime selector, terminal matrix, Koide determinant map, and
canonical period convention are taken from the byte-preserving current source
bundle supplied with this continuation.  Its embedded files verify against
their recorded SHA-256 values.  The original complex-threefold construction
is attributed to Levent Alpöge with the disclosed Claude/Fable contribution;
the separate Erdős--Straus research archive is attributed in its release to
The Clankers.  No theorem below uses the unverified global identification of
the constructed threefold with $S^6$.  Only the displayed arithmetic matrices
and the ambient period convention are used.

\section{Imported prime selector and determinant dictionaries}

Let $p\equiv1\pmod{12}$ be prime.  For
\begin{equation}\label{eq:R-a-prime}
 R_j=4j+3,\qquad a_j=\frac{p+R_j}{4},
 \qquad 0\le j<\frac{p-1}{4},
\end{equation}
the two source selector sets are
\begin{align}
 \mathcal W_0(p,j)
 &=\{u\in\Z_{>0}:u\mid a_j^2,\ R_j\mid4u+1\},\label{eq:E-source}\\
 \mathcal W_1(p,j)
 &=\{u\in\Z_{>0}:u\mid a_j^2,\ R_j\mid4u+p\}.\label{eq:M-source}
\end{align}
They are the exterior and middle channels.  With $\eps\in\{0,1\}$ and
$u\in\mathcal W_\eps(p,j)$, put
\begin{equation}\label{eq:source-T}
 T=\frac{4u+p^\eps}{R_j},\qquad
 (\ell_0,\ell_1,\ell_2)
 =\left(\frac{T+1}{4},u-jT,1\right).
\end{equation}
Then
\begin{equation}\label{eq:source-lambda}
 12\ell_0-4\ell_1-3\ell_2=p^\eps
\end{equation}
and
\begin{equation}\label{eq:source-shear}
 \mat{1}{-j}{0}{1}
 \mat{R_j}{-u}{4}{-T}
 =
 \mat{3}{-\ell_1}{4}{1-4\ell_0}.
\end{equation}
Thus the selector presentation and the terminal presentation differ by an
explicit element of $\SL_2(\Z)$.

The rank-one determinant model used below is
\begin{equation}\label{eq:iota}
 \iota(H,X,Y)=\mat{H+Y}{X}{X}{H-Y},\qquad
 \det\iota(H,X,Y)=H^2-X^2-Y^2.
\end{equation}
The source orthogonal Koide map $\mathcal F:\R^3\to\R^3$ satisfies
\begin{equation}\label{eq:koide-orthogonal}
 Q_K(q)=\frac23(q_1+q_2+q_3)^2-(q_1^2+q_2^2+q_3^2)
 =H^2-X^2-Y^2,
\end{equation}
where
\begin{equation}\label{eq:F}
 \begin{pmatrix}H\\X\\Y\end{pmatrix}
 =
 \begin{pmatrix}
 1/\sqrt3&1/\sqrt3&1/\sqrt3\\
 1/\sqrt2&0&-1/\sqrt2\\
 1/\sqrt6&-2/\sqrt6&1/\sqrt6
 \end{pmatrix}
 \begin{pmatrix}q_1\\q_2\\q_3\end{pmatrix}.
\end{equation}

Two coefficient dictionaries are retained separately.  The native symmetric
coefficient matrix is
\begin{equation}\label{eq:native-G}
 G(t,m,b)=\mat{2t}{2m}{2m}{b/3},
\end{equation}
with
\begin{equation}\label{eq:L-native}
 \mathcal L(t,m,b)=\left(t+\frac b6,2m,t-\frac b6\right)=(H,X,Y),
\end{equation}
so
\begin{equation}\label{eq:native-disc}
 \det G=\frac23(tb-6m^2).
\end{equation}
The raw-period dictionary is
\begin{equation}\label{eq:raw-dict}
 \mathcal P_K(H,X,Y)=\left(H+Y,\frac X{\sqrt6},H-Y\right)
 =(T_{\rm raw},M_{\rm raw},B_{\rm raw}),
\end{equation}
and
\begin{equation}\label{eq:raw-disc}
 T_{\rm raw}B_{\rm raw}-6M_{\rm raw}^2=H^2-X^2-Y^2.
\end{equation}
Their exact connecting map is
\begin{equation}\label{eq:D-map}
 \mathcal D(t,m,b)=\left(2t,\sqrt{\frac23}\,m,\frac b3\right),
 \qquad \mathcal P_K=\mathcal D\circ\mathcal L^{-1}.
\end{equation}

\section{Complete divisor channels for a coprime odd shell}

Fix integers
\begin{equation}\label{eq:general-shell-hyp}
 n\equiv1\pmod4,\qquad R\equiv3\pmod4,
 \qquad 0<R<n,\qquad \gcd(n,R)=1,
\end{equation}
and put
\begin{equation}\label{eq:aN}
 a=\frac{n+R}{4},\qquad N=na.
\end{equation}
Then
\begin{equation}\label{eq:gcd-general}
 \gcd(n,a)=\gcd(R,a)=1.
\end{equation}
The fixed-first-denominator equation
\begin{equation}\label{eq:fixed-shell-ES}
 \frac4n=\frac1a+\frac1y+\frac1z
\end{equation}
is equivalent to
\begin{equation}\label{eq:residual-equation}
 Ryz=N(y+z)
\end{equation}
and to
\begin{equation}\label{eq:shifted-factorization}
 (Ry-N)(Rz-N)=N^2.
\end{equation}

\begin{definition}[oriented divisor channel]\label{def:channel}
For $\delta\mid n^2$, define
\begin{equation}\label{eq:channel-set}
 \mathscr C_\delta(n,R)=
 \left\{(u,v)\in\Z_{>0}^2:
 uv=a^2,\quad R\mid4u+\delta\right\}.
\end{equation}
For $(\delta,u,v)$ in this set, define
\begin{equation}\label{eq:d-e-oriented}
 d=\frac{n^2}{\delta}u,\qquad e=\delta v
\end{equation}
and
\begin{equation}\label{eq:y-z-general}
 y=\frac{N+d}{R},\qquad z=\frac{N+e}{R}.
\end{equation}
The orientation in \eqref{eq:d-e-oriented} is fixed so that, when $n=p$ is
prime, $\delta=1$ is exactly the source exterior channel
$d=p^2u$.
\end{definition}

\begin{theorem}[complete oriented channel decomposition]\label{thm:complete-channel}
The map
\begin{equation}\label{eq:channel-map}
 \bigsqcup_{\delta\mid n^2}\mathscr C_\delta(n,R)
 \longrightarrow
 \left\{(y,z)\in\Z_{>0}^2:Ryz=N(y+z)\right\}
\end{equation}
given by \eqref{eq:d-e-oriented}--\eqref{eq:y-z-general} is a bijection.
Its inverse is
\begin{align}
 d&=Ry-N,&e&=Rz-N,\label{eq:inverse-de}\\
 \kappa&=\gcd(d,n^2),&
 \delta&=\frac{n^2}{\kappa},\label{eq:inverse-delta}\\
 u&=\frac d\kappa,&v&=\frac e\delta.\label{eq:inverse-uv}
\end{align}
Under exchange of the final denominators, the channel coordinates transform
as
\begin{equation}\label{eq:complement-involution}
 (\delta,u,v)\longmapsto
 \left(\frac{n^2}{\delta},v,u\right).
\end{equation}
The involution has no fixed point.  Therefore the number of unordered
solutions in this fixed shell is
\begin{equation}\label{eq:unordered-half}
 \frac12\sum_{\delta\mid n^2}|\mathscr C_\delta(n,R)|.
\end{equation}
\end{theorem}

\begin{proof}
Suppose first that $(\delta,u,v)$ satisfies \eqref{eq:channel-set}.  Then
\[
 de=\frac{n^2}{\delta}u\cdot\delta v=n^2a^2=N^2.
\]
To verify the first integrality in \eqref{eq:y-z-general}, multiply its
numerator by the unit $\delta$ modulo $R$:
\[
 \delta(N+d)=\delta na+n^2u
 =n(\delta a+nu).
\]
Because $n\equiv4a\pmod R$, the final parenthesis is congruent to
$a(\delta+4u)$, hence to zero.  Since $\delta$ is a unit modulo $R$,
$R\mid N+d$.  For the second numerator,
\[
 N+e=na+\delta v\equiv4a^2+\delta v
 =v(4u+\delta)\equiv0\pmod R.
\]
Both numbers are positive, and \eqref{eq:shifted-factorization} proves the
reciprocal equation.

Conversely, let $(y,z)$ satisfy \eqref{eq:residual-equation}.  Division by
$z$ and by $y$ shows $d=Ry-N>0$ and $e=Rz-N>0$, while
\eqref{eq:shifted-factorization} gives $de=n^2a^2$.  Since $\gcd(n,a)=1$,
the $n$-primary and $a$-primary parts split uniquely.  In particular
$\kappa=\gcd(d,n^2)$, $\delta=n^2/\kappa$, $u=d/\kappa$, and $v=e/\delta$
are positive integers with $uv=a^2$.  Finally,
\[
 0\equiv N+e=na+\delta v\equiv4a^2+\delta v
 =v(4u+\delta)\pmod R.
\]
The integer $v$ is a unit modulo $R$, so $R\mid4u+\delta$.  The two
constructions are literal inverses.

Exchanging $y,z$ exchanges $d,e$ and gives \eqref{eq:complement-involution}.
A fixed point would require $\delta=n$ and $u=v=a$.  Its channel congruence
would give
\[
 R\mid4a+n=2n+R,
\]
so $R\mid2n$.  This contradicts $R>1$, $R$ odd, and $\gcd(R,n)=1$.
\end{proof}

\begin{corollary}[prime specialization]\label{cor:prime-specialization}
For $n=p$ prime, the three oriented channel orders are
\begin{equation}\label{eq:prime-deltas}
 \delta=1,\quad p,\quad p^2.
\end{equation}
The first is the source exterior channel, the second is the source middle
channel, and the third is the reversed exterior orientation.  Hence
\begin{equation}\label{eq:prime-count}
 \#\{\text{unordered solutions in the shell}\}
 =|\mathcal W_0(p,j)|+\frac12|\mathcal W_1(p,j)|.
\end{equation}
\end{corollary}

\begin{example}[the orientation correction]\label{ex:n9}
Take
\[
 n=9,\qquad R=7,\qquad a=4,\qquad \delta=3,
 \qquad u=1,\qquad v=16.
\]
Then
\[
 d=\frac{81}{3}\cdot1=27,\qquad e=3\cdot16=48,
\]
and
\[
 y=\frac{36+27}{7}=9,\qquad z=\frac{36+48}{7}=12.
\]
Thus
\[
 \frac49=\frac14+\frac19+\frac1{12}.
\]
The inverse gives $\gcd(d,81)=27$ and therefore
$\delta=81/27=3$.  Replacing this by $\gcd(d,81)$ would reverse the channel
order and would not preserve the source exterior orientation.
\end{example}

\section{Terminal matrices and the divisor-order spectrum}

For a channel state $(\delta,u,v)$, put
\begin{equation}\label{eq:Tdelta}
 T_\delta=\frac{4u+\delta}{R},\qquad
 j=\frac{R-3}{4},
\end{equation}
and define
\begin{equation}\label{eq:sigdelta}
 \sigma_\delta=
 \begin{cases}
 1,&\delta\equiv1\pmod4,\\
 -1,&\delta\equiv3\pmod4.
 \end{cases}
\end{equation}
Because $RT_\delta\equiv\delta\pmod4$ and $R\equiv-1\pmod4$,
\begin{equation}\label{eq:T-sigma}
 T_\delta+\sigma_\delta\equiv0\pmod4.
\end{equation}
Set
\begin{equation}\label{eq:ell-general}
 \ell_2=\sigma_\delta,
 \qquad
 \ell_0=\frac{T_\delta+\sigma_\delta}{4},
 \qquad
 \ell_1=u-jT_\delta.
\end{equation}

\begin{theorem}[complete terminal matrix of a divisor channel]\label{thm:terminal-general}
The coordinates \eqref{eq:ell-general} are integral and satisfy
\begin{equation}\label{eq:lambda-delta}
 12\ell_0-4\ell_1-3\ell_2=\delta.
\end{equation}
Moreover,
\begin{equation}\label{eq:general-row-shear}
 \mat{1}{-j}{0}{1}
 \mat{R}{-u}{4}{-T_\delta}
 =
 \underbrace{\mat{3}{-\ell_1}{4}{\ell_2-4\ell_0}}
 _{\mathsf R_\delta}
 =\mat{3}{-\ell_1}{4}{-T_\delta}.
\end{equation}
Consequently
\begin{equation}\label{eq:terminal-snf}
 \det\mathsf R_\delta=-\delta,
 \qquad
 \SNF(\mathsf R_\delta)=\diag(1,\delta).
\end{equation}
Finally,
\begin{equation}\label{eq:ell1mod3}
 \ell_1\equiv-\delta\pmod3.
\end{equation}
Thus the source admissibility $3\nmid\ell_1$ holds exactly when
$3\nmid\delta$, and $\ell_2$ is always odd.
\end{theorem}

\begin{proof}
Only the displayed identities are needed.  Equation \eqref{eq:T-sigma}
proves $\ell_0\in\Z$.  Since $R=4j+3$,
\begin{align*}
 12\ell_0-4\ell_1-3\ell_2
 &=3(T_\delta+\sigma_\delta)-4(u-jT_\delta)
   -3\sigma_\delta\\
 &=(4j+3)T_\delta-4u=RT_\delta-4u=\delta.
\end{align*}
Left multiplication in \eqref{eq:general-row-shear} gives first row
$(R-4j,-u+jT_\delta)=(3,-\ell_1)$; its second row is unchanged.  Also
$\ell_2-4\ell_0=-T_\delta$.  The entries $3,4$ have gcd one, so the first
Smith invariant is one; the determinant gives the second.  Reducing
$\delta=3T_\delta-4\ell_1$ modulo three proves \eqref{eq:ell1mod3}.
\end{proof}

\begin{corollary}[complementary terminal orders]\label{cor:terminal-complement}
The two orientations of one unordered reciprocal solution have terminal
orders
\begin{equation}\label{eq:paired-orders}
 \left\{\delta,\frac{n^2}{\delta}\right\},
 \qquad
 \delta\cdot\frac{n^2}{\delta}=n^2.
\end{equation}
The unique self-complementary order is $\delta=n$, but the oriented channel
involution has no fixed state.
\end{corollary}

\begin{remark}[scope of the terminal homology package]
The canonical workbench contains a terminal chain calculation which, after
its displayed geometric inputs are granted, assigns
\[
 (H_0,\ldots,H_6)
 \cong
 (\Z,\Z/\delta,\Z/\delta,\Z/\delta,\Z/\delta,0,\Z)
\]
to a nonzero terminal value $\delta$.  In particular $\delta=1$ is the
homology-sphere package.  The present theorem supplies its exact arithmetic
input matrix; it does not independently reprove that every such matrix is
realized by the global complex-threefold gluing.
\end{remark}

\section{Signed divisor boxes and exact Fourier counts}

Factor
\begin{equation}\label{eq:a-factor}
 a=\prod_{q^e\parallel a}q^e.
\end{equation}
Every divisor $u\mid a^2$ has a unique signed exponent vector
\begin{equation}\label{eq:signed-exponents}
 u=a\prod_{q^e\parallel a}q^{\beta_q},
 \qquad -e\le\beta_q\le e.
\end{equation}
Since $4a\equiv n\pmod R$, the channel condition becomes
\begin{equation}\label{eq:signed-channel}
 \prod_{q^e\parallel a}q^{\beta_q}
 \equiv-\delta n^{-1}\pmod R.
\end{equation}
All terms are units modulo $R$ by \eqref{eq:gcd-general}.

Let
\[
 U_R=(\Z/R\Z)^\times,
 \qquad \widehat U_R=\operatorname{Hom}(U_R,\C^\times),
\]
and define
\begin{equation}\label{eq:Dchi}
 D_a(\chi)=
 \prod_{q^e\parallel a}\left(\sum_{\beta=-e}^{e}\chi(q)^\beta\right).
\end{equation}

\begin{theorem}[exact channel Fourier formula]\label{thm:fourier}
For every $\delta\mid n^2$,
\begin{equation}\label{eq:Cdelta-fourier}
 C_\delta(n,R):=|\mathscr C_\delta(n,R)|
 =\frac1{\varphi(R)}
 \sum_{\chi\in\widehat U_R}
 \overline{\chi(-\delta n^{-1})}\,D_a(\chi).
\end{equation}
If $n=\prod_{q^f\parallel n}q^f$, then the number of unordered reciprocal
solutions in the shell is
\begin{equation}\label{eq:total-fourier}
 \boxed{
 \frac1{2\varphi(R)}
 \sum_{\chi\in\widehat U_R}
 \overline{\chi(-n^{-1})}\,D_a(\chi)
 \prod_{q^f\parallel n}
 \left(\sum_{j=0}^{2f}\overline{\chi(q)}^{\,j}\right).}
\end{equation}
The principal-character contribution is exactly
\begin{equation}\label{eq:principal-term}
 \frac{\tau(a^2)\tau(n^2)}{2\varphi(R)}.
\end{equation}
\end{theorem}

\begin{proof}
Character orthogonality on $U_R$ gives, for $x,c\in U_R$,
\[
 \1_{x=c}=\frac1{\varphi(R)}
 \sum_{\chi\in\widehat U_R}\chi(x)\overline{\chi(c)}.
\]
Sum this identity over the signed box \eqref{eq:signed-exponents}, with
$c=-\delta n^{-1}$.  The independent exponent sums factor and give
\eqref{eq:Cdelta-fourier}.  Summing over $\delta\mid n^2$ gives
\[
 \sum_{\delta\mid n^2}\overline{\chi(\delta)}
 =\prod_{q^f\parallel n}\sum_{j=0}^{2f}\overline{\chi(q)}^{\,j}.
\]
The fixed-point-free involution in \cref{thm:complete-channel} divides the
oriented total by two.  For the principal character,
$D_a(1)=\tau(a^2)$ and the divisor sum is $\tau(n^2)$.
\end{proof}

\section{Primitive rays and the exact Lorentz channel action}

\begin{lemma}[canonical primitive factor coordinates]\label{lem:factor-ray}
For $u,v,a>0$ with $uv=a^2$, there are unique positive integers $r,s,t$ such
that
\begin{equation}\label{eq:rst}
 u=r^2s,\qquad a=rst,\qquad v=st^2,
 \qquad \gcd(r,t)=1.
\end{equation}
\end{lemma}

\begin{proof}
Let $g=\gcd(u,a)$ and set
\[
 r=\frac ug,\qquad t=\frac ag,\qquad s=\frac gr.
\]
Prime valuations in $u\mid a^2$ show $r\mid g$, so $s\in\Z_{>0}$.
The displayed identities and $\gcd(r,t)=1$ follow immediately.  Conversely,
$g=rs$, so the construction recovers the same triple and proves uniqueness.
\end{proof}

Since $n=4rst-R$, one has the exact identity
\begin{equation}\label{eq:projective-identity}
 nr+\delta t=t(4r^2s+\delta)-Rr.
\end{equation}
Because $r,t$ are units modulo $R$,
\begin{equation}\label{eq:projective-incidence}
 \boxed{R\mid4r^2s+\delta
 \quad\Longleftrightarrow\quad R\mid nr+\delta t.}
\end{equation}
Thus the channel is the inverse image of $[1:-1]$ under the exact projective
map
\begin{equation}\label{eq:projective-map}
 \Phi_{n,\delta}:\Pone(\Z/R\Z)\longrightarrow\Pone(\Z/R\Z),
 \qquad [r:t]\longmapsto[nr:\delta t],
\end{equation}
restricted to the affine chart in which both coordinates are units.

The rank-one matrix of the ray is
\begin{equation}\label{eq:Gamma-ray}
 \Gamma_{r,s,t}
 =s\begin{pmatrix}r\\t\end{pmatrix}
   \begin{pmatrix}r&t\end{pmatrix}
 =\mat{u}{a}{a}{v}.
\end{equation}
The shifted shell matrix is
\begin{equation}\label{eq:Sigma-channel}
 \Sigma_{n,\delta}
 =\mat{(n^2/\delta)u}{na}{na}{\delta v}.
\end{equation}
Let
\begin{equation}\label{eq:Mndelta}
 M_{n,\delta}=\diag\left(\frac n{\sqrt\delta},\sqrt\delta\right).
\end{equation}
Then
\begin{equation}\label{eq:congruence-action}
 \boxed{\Sigma_{n,\delta}
 =M_{n,\delta}\Gamma_{r,s,t}M_{n,\delta}^{\mathsf T}.}
\end{equation}

\begin{theorem}[integral Lorentz similitude]\label{thm:lorentz}
Under \eqref{eq:iota}, congruence by $M_{n,\delta}$ is the integral linear map
\begin{equation}\label{eq:rho}
 \rho_{n,\delta}=
 \begin{pmatrix}
 \dfrac{n^2/\delta+\delta}{2}&0&
 \dfrac{n^2/\delta-\delta}{2}\\[1ex]
 0&n&0\\[1ex]
 \dfrac{n^2/\delta-\delta}{2}&0&
 \dfrac{n^2/\delta+\delta}{2}
 \end{pmatrix}.
\end{equation}
Its entries are integers and, with $J_L=\diag(1,-1,-1)$,
\begin{equation}\label{eq:lorentz-similitude}
 \rho_{n,\delta}^{\mathsf T}J_L\rho_{n,\delta}=n^2J_L,
 \qquad
 \det\rho_{n,\delta}=n^3.
\end{equation}
Writing
\begin{equation}\label{eq:rapidity}
 \xi_\delta=\log\frac n\delta,
\end{equation}
one has
\begin{equation}\label{eq:boost}
 \rho_{n,\delta}=n
 \begin{pmatrix}
 \cosh\xi_\delta&0&\sinh\xi_\delta\\
 0&1&0\\
 \sinh\xi_\delta&0&\cosh\xi_\delta
 \end{pmatrix}.
\end{equation}
Complementary channels have opposite rapidities, and the middle channel
$\delta=n$ is the pure dilation $nI_3$.
\end{theorem}

\begin{proof}
For a diagonal matrix $\diag(A,B)$, congruence sends
$H+Y$ to $A^2(H+Y)$, $H-Y$ to $B^2(H-Y)$, and $X$ to $ABX$.
Substitution of $A=n/\sqrt\delta$, $B=\sqrt\delta$ gives
\eqref{eq:rho}.  Since $n$ and $\delta$ are odd, both half-sums are integers.
Direct multiplication gives \eqref{eq:lorentz-similitude}.  Finally,
\[
 n\cosh\log(n/\delta)=\frac{n^2/\delta+\delta}{2},\qquad
 n\sinh\log(n/\delta)=\frac{n^2/\delta-\delta}{2}.
\]
Replacing $\delta$ by $n^2/\delta$ negates the logarithm.
\end{proof}

The affine ray slope $\theta=t/r$ transforms under $M_{n,\delta}$ as
\begin{equation}\label{eq:slope-transform}
 \theta\longmapsto\theta_\delta=\frac\delta n\theta.
\end{equation}
The channel incidence \eqref{eq:projective-incidence} is exactly
\begin{equation}\label{eq:universal-middle-incidence}
 \boxed{\theta_\delta\equiv-1\pmod R.}
\end{equation}
Thus every divisor channel is carried by its displayed Lorentz morphism to
the same affine middle incidence.

\section{Prime saturation on every primitive null ray}

Fix a primitive positive pair $(r,t)$ and an arithmetic progression
\begin{equation}\label{eq:M-c}
 4\mid M,\qquad \gcd(c,M)=1,\qquad c\equiv1\pmod4.
\end{equation}
Put
\begin{equation}\label{eq:A-g}
 A=4rt,
 \qquad g=\gcd(A,M),
 \qquad M'=\frac Mg.
\end{equation}
Choose a prime $R$ such that
\begin{equation}\label{eq:R-choice}
 R\equiv-c\pmod g,
 \qquad R\equiv3\pmod4,
 \qquad \gcd(R,rtM)=1.
\end{equation}
The first congruence already implies the second because $4\mid g$ and
$c\equiv1\pmod4$.  Such primes $R$ exist by Dirichlet's theorem.

Define $s_R\pmod R$ and $s_M\pmod{M'}$ by
\begin{align}
 4r^2s_R&\equiv-1\pmod R,\label{eq:sR}\\
 \frac Ags_M&\equiv\frac{c+R}{g}\pmod{M'}.
 \label{eq:sM}
\end{align}
All displayed coefficients are units in their respective rings, and
$\gcd(R,M')=1$.  Let $s_0$ be the unique class modulo $RM'$ satisfying both.
Put
\begin{equation}\label{eq:p-progression}
 p_0=As_0-R,
 \qquad
 Q=ARM'=MR\frac Ag.
\end{equation}

\begin{theorem}[prime saturation of each primitive ray]\label{thm:prime-saturation}
The residue $p_0\pmod Q$ is reduced.  Every sufficiently large prime
\begin{equation}\label{eq:p-AP}
 p\equiv p_0\pmod Q
\end{equation}
satisfies $p\equiv c\pmod M$ and has an exterior Erdős--Straus certificate on
the exact primitive ray $[r:t]$.  Explicitly, write
\begin{equation}\label{eq:s-from-p}
 s=\frac{p+R}{4rt},
 \qquad
 a=rst,
 \qquad
 u=r^2s,
 \qquad
 v=st^2.
\end{equation}
Then
\begin{equation}\label{eq:E-ray-incidence}
 R\mid pr+t,
 \qquad
 R\mid4u+1,
\end{equation}
and the reciprocal denominators are
\begin{equation}\label{eq:E-denominators}
 \boxed{
 x=rst,
 \qquad
 y=prs\,\frac{pr+t}{R},
 \qquad
 z=st\,\frac{pr+t}{R}.}
\end{equation}
Consequently
\begin{equation}\label{eq:PNT-ray}
 \#\{p\le X:p\equiv p_0\pmod Q\}
 \sim\frac{\Li(X)}{\varphi(Q)}.
\end{equation}
Within the primes $p\equiv c\pmod M$, the constructed subfamily has relative
density
\begin{equation}\label{eq:relative-density}
 \frac{\varphi(M)}{\varphi(Q)}>0.
\end{equation}
\end{theorem}

\begin{proof}
The congruences \eqref{eq:sR}--\eqref{eq:sM} imply
$As_0-R\equiv c\pmod M$ and $4r^2s_0+1\equiv0\pmod R$.
The modulus in \eqref{eq:p-progression} is exactly the change in $p=As-R$
when $s$ is increased by $RM'$.  To prove $\gcd(p_0,Q)=1$, note
\[
 p_0\equiv-R\pmod A,
 \qquad
 p_0\equiv As_R\equiv-tr^{-1}\pmod R,
 \qquad
 p_0\equiv c\pmod{M'}.
\]
Each right side is a unit in the indicated modulus.  Dirichlet's theorem, or
the prime number theorem in arithmetic progressions, gives
\eqref{eq:PNT-ray}.

For such a prime, \eqref{eq:s-from-p} is integral and gives $4a=p+R$ and
$uv=a^2$.  The first congruence in \eqref{eq:E-ray-incidence} follows from
\[
 pr+t=t(4r^2s+1)-Rr.
\]
The second is equivalent because $t$ is a unit modulo $R$.  Formula
\eqref{eq:E-denominators} is the source factor-coordinate reconstruction;
direct substitution also gives
$4xyz=p(xy+xz+yz)$.  Dividing the two prime-counting asymptotics for moduli
$Q$ and $M$ gives \eqref{eq:relative-density}.
\end{proof}

\begin{corollary}[all $48$ reduced $1\bmod12$ classes modulo $840$]
For every primitive positive ray $[r:t]$ and every one of the $48$ reduced
classes
\begin{equation}\label{eq:48classes}
 c\pmod{840},\qquad c\equiv1\pmod{12},
\end{equation}
there is a reduced prime progression consisting entirely, apart from a finite
initial segment, of exterior certificates on that ray.
\end{corollary}

\section{A single Koide ray captures density one of every progression}

The preceding theorem lets $R$ depend on the ray and progression.  On the
diagonal ray $[1:1]$, a stronger fixed description is available.

Let $p\equiv1\pmod4$ be prime and let $q\equiv3\pmod4$ be a prime divisor of
$p+1$.  Put
\begin{equation}\label{eq:diagonal-sk}
 s=\frac{p+q}{4},\qquad k=\frac{p+1}{q}.
\end{equation}
Then $u=a=v=s$ and
\begin{equation}\label{eq:diagonal-identity}
 \boxed{
 \frac4p=\frac1s+\frac1{sk}+\frac1{psk}.}
\end{equation}
Indeed,
\[
 \frac1s+\frac1{sk}+\frac1{psk}
 =\frac{pk+p+1}{psk}
 =\frac{(p+1)(p+q)/q}{p(p+q)(p+1)/(4q)}
 =\frac4p.
\]
Since $qk=p+1\equiv2\pmod4$ and $q\equiv-1\pmod4$, one has
$k\equiv2\pmod4$.  This is the exterior channel with
\begin{equation}\label{eq:diagonal-terminal}
 R=q,
 \quad j=\frac{q-3}{4},
 \quad T=k+1,
 \quad
 (\ell_0,\ell_1,\ell_2)
 =\left(\frac{k+2}{4},\frac{3k+2}{4},1\right),
\end{equation}
and terminal determinant $-1$.

Fix $M,c$ as in \eqref{eq:M-c}.  For a finite set $\mathcal Q$ of primes
$q\equiv3\pmod4$ with $q\nmid M$, let $\mathscr P_{\mathcal Q}(M,c)$ denote
the primes $p\equiv c\pmod M$ for which $q\mid p+1$ for at least one
$q\in\mathcal Q$.

\begin{theorem}[finite sieve density and density-one diagonal ray]
\label{thm:density-one-ray}
The relative density of $\mathscr P_{\mathcal Q}(M,c)$ inside the primes
$p\equiv c\pmod M$ exists and equals
\begin{equation}\label{eq:finite-union-density}
 \boxed{
 1-\prod_{q\in\mathcal Q}\frac{q-2}{q-1}.}
\end{equation}
As $\mathcal Q$ increases through all primes $q\equiv3\pmod4$, $q\nmid M$,
the relative density tends to one.  Consequently the exterior certificates
on the single primitive Koide ray
\begin{equation}\label{eq:diagonal-ray}
 [r:t]=[1:1],
 \qquad
 \Gamma_s=s\mat{1}{1}{1}{1},
\end{equation}
capture relative density one of the primes in every reduced progression
$c\pmod M$ with $c\equiv1\pmod4$.
\end{theorem}

\begin{proof}
Put $Q_{\mathcal Q}=\prod_{q\in\mathcal Q}q$.  Since every $q$ is coprime to
$M$, the Chinese remainder theorem identifies the reduced residue classes
modulo $MQ_{\mathcal Q}$ above $c\pmod M$ with the product of the
$q-1$ unit classes modulo each $q$.  Avoiding $p\equiv-1\pmod q$ leaves
$q-2$ classes.  The prime number theorem in arithmetic progressions therefore
gives relative density
\[
 \prod_{q\in\mathcal Q}\frac{q-2}{q-1}
\]
for the complement, proving \eqref{eq:finite-union-density}.

The sum of reciprocals of primes $q\equiv3\pmod4$ diverges.  Hence
\[
 \sum_{q\equiv3(4)}\frac1{q-1}=\infty,
 \qquad
 \prod_{q\in\mathcal Q}
 \left(1-\frac1{q-1}\right)\longrightarrow0.
\]
Every prime counted by the union has the explicit certificate
\eqref{eq:diagonal-identity}; taking larger finite sets shows that the
relative lower density of the union is at least every quantity in
\eqref{eq:finite-union-density}, hence is one.
\end{proof}

\begin{corollary}[a fixed $10\%$ subprogression modulo $9240$]
Take $M=840$ and $q=11$.  In each reduced class $c\equiv1\pmod{12}$ modulo
$840$, the simultaneous congruences
\begin{equation}\label{eq:q11-CRT}
 p\equiv c\pmod{840},
 \qquad p\equiv-1\pmod{11}
\end{equation}
define one reduced class modulo $9240$.  Its primes have the diagonal-ray
certificate \eqref{eq:diagonal-identity} and form relative density
$1/(11-1)=1/10$ inside the class $c\pmod{840}$.
\end{corollary}

\begin{example}
For $p=109$ and $q=11$,
\[
 s=30,\qquad k=10,
\]
so
\begin{equation}\label{eq:p109}
 \frac4{109}=\frac1{30}+\frac1{300}+\frac1{32700}.
\end{equation}
The terminal triple is $(\ell_0,\ell_1,\ell_2)=(3,8,1)$ and
\[
 \mathsf R=\mat{3}{-8}{4}{-11},\qquad \det\mathsf R=-1.
\]
\end{example}

\begin{remark}[the unresolved quantifier]
The theorem proves
\[
 \forall\,[r:t]\ \forall\,c\pmod M\quad
 \text{infinitely many primes in the class have a certificate on the ray},
\]
and, on $[1:1]$, a density-one statement.  The Erdős--Straus conjecture asks
instead
\[
 \forall\,p\quad\exists\,[r:t],R\quad\text{a certificate for that particular
 prime}.
\]
A density-zero exceptional set is not removed by \cref{thm:density-one-ray}.
\end{remark}



\section{The universal diagonal divisor identity}

Retain the complete coprime odd shell \eqref{eq:general-shell-hyp} and restrict
to the primitive ray $[r:t]=[1:1]$.  Then $u=a=v=s$ with
$s=(n+R)/4$.

\begin{proposition}[all diagonal divisor channels]
\label{prop:all-diagonal-channels}
For $\delta\mid n^2$, the diagonal ray supports the oriented channel of order
$\delta$ exactly when
\begin{equation}\label{eq:general-diagonal-incidence}
 R\mid n+\delta.
\end{equation}
Put
\begin{equation}\label{eq:general-diagonal-k}
 k_\delta=\frac{n+\delta}{R}.
\end{equation}
Then the corresponding reciprocal decomposition is
\begin{equation}\label{eq:general-diagonal-decomposition}
 \boxed{
 \frac4n
 =\frac1s+\frac1{s k_\delta}
  +\frac1{s n k_\delta/\delta},
 \qquad s=\frac{n+R}{4}.}
\end{equation}
The third displayed denominator is an integer.  The terminal coordinates are
\begin{equation}\label{eq:general-diagonal-terminal}
 T_\delta=k_\delta+1,
 \qquad
 \ell_2=\sigma_\delta,
 \qquad
 \ell_0=\frac{k_\delta+1+\sigma_\delta}{4},
 \qquad
 \ell_1=\frac{3k_\delta+3-\delta}{4},
\end{equation}
and their terminal determinant is $-\delta$.

Under the complementary orientation
$\delta'=n^2/\delta$, one has
\begin{equation}\label{eq:general-diagonal-complement-k}
 k_{\delta'}=\frac n\delta k_\delta,
\end{equation}
and the final two denominators in
\eqref{eq:general-diagonal-decomposition} are exchanged exactly.
\end{proposition}

\begin{proof}
On $[1:1]$, equation \eqref{eq:projective-incidence} is
$R\mid n+\delta$, proving the first assertion.  The channel reconstruction
\eqref{eq:y-z-general} gives
\[
 y=\frac{ns+(n^2/\delta)s}{R}=\frac{snk_\delta}{\delta},
 \qquad
 z=\frac{ns+\delta s}{R}=sk_\delta.
\]
Their integrality is already part of the complete channel theorem; direct
addition gives
\[
 \frac1s+\frac1{sk_\delta}+\frac{\delta}{snk_\delta}
 =\frac{nk_\delta+n+\delta}{snk_\delta}
 =\frac{k_\delta(n+R)}{snk_\delta}=\frac4n.
\]
Because $T_\delta=(4s+\delta)/R$, one has
$T_\delta=k_\delta+1$.  Substitution in \eqref{eq:ell-general} gives
\eqref{eq:general-diagonal-terminal}; its integrality and determinant are
those of \cref{thm:terminal-general}.  Finally,
\[
 k_{n^2/\delta}
 =\frac{n+n^2/\delta}{R}
 =\frac n\delta\frac{n+\delta}{R},
\]
and substitution swaps $sk_\delta$ with $snk_\delta/\delta$.
\end{proof}

\section{Exact multiplicity on the diagonal Koide ray}

Let $p\equiv1\pmod4$ be prime and put
\begin{equation}\label{eq:diag-m}
 m_p=\frac{p+1}{2}.
\end{equation}
Thus $m_p$ is odd.  Let $\chi_4$ denote the nontrivial real Dirichlet
character modulo four on odd integers.

\begin{theorem}[complete diagonal-ray certificate count]
\label{thm:diag-count}
The map
\begin{equation}\label{eq:diag-R-map}
 R\longmapsto
 \left(
  x_R,y_R,z_R
 \right)
 =\left(
  \frac{p+R}{4},
  \frac{p+R}{4}\frac{p+1}{R},
  p\frac{p+R}{4}\frac{p+1}{R}
 \right)
\end{equation}
is a bijection from
\begin{equation}\label{eq:diag-R-set}
 \mathscr R_{\mathrm{diag}}(p)
 =\{R>0:R\mid m_p,\ R\equiv3\pmod4\}
\end{equation}
onto the exterior Erdős--Straus selector states on the exact primitive ray
$[r:t]=[1:1]$.  Every displayed triple is ordered and satisfies
\begin{equation}\label{eq:diag-unit-fraction}
 \frac4p=\frac1{x_R}+\frac1{y_R}+\frac1{z_R}.
\end{equation}
The number of these states is
\begin{equation}\label{eq:diag-count-character}
 \boxed{
 \nu_{\mathrm{diag}}(p)
 =\frac12\left(
   \tau(m_p)-\sum_{d\mid m_p}\chi_4(d)
  \right).}
\end{equation}
If
\begin{equation}\label{eq:diag-factorization}
 m_p=
 \prod_{q\equiv1(4)}q^{e_q}
 \prod_{r\equiv3(4)}r^{f_r},
\end{equation}
then
\begin{equation}\label{eq:diag-count-factorized}
 \boxed{
 \nu_{\mathrm{diag}}(p)=\frac12\left[
 \prod_{q\equiv1(4)}(e_q+1)
 \prod_{r\equiv3(4)}(f_r+1)
 -
 \prod_{q\equiv1(4)}(e_q+1)
 \prod_{r\equiv3(4)}\epsilon(f_r)
 \right],}
\end{equation}
where $\epsilon(f)=1$ for even $f$ and $\epsilon(f)=0$ for odd $f$.
Consequently
\begin{equation}\label{eq:diag-exceptional-exact}
 \boxed{
 \nu_{\mathrm{diag}}(p)=0
 \quad\Longleftrightarrow\quad
 \text{every prime divisor of }(p+1)/2\text{ is }1\pmod4.}
\end{equation}
\end{theorem}

\begin{proof}
On the ray $[1:1]$, the canonical factor coordinates are
$u=a=v=s$.  The exterior projective incidence
\eqref{eq:projective-incidence} becomes $R\mid p+1$, while the shell equation
$4a=p+R$ gives $s=(p+R)/4$.  Because $R$ is odd, $R\mid p+1$ is equivalent to
$R\mid m_p$.  Conversely, every divisor in \eqref{eq:diag-R-set} satisfies
$3\le R\le m_p\le p-2$ and defines one allowed shell marker.  With
$k_R=(p+1)/R$, the source exterior reconstruction is exactly
$(x_R,y_R,z_R)=(s,sk_R,psk_R)$, proving the bijection and
\eqref{eq:diag-unit-fraction}.  Since $k_R\ge2$, the triple is ordered.
Distinct values of $R$ give distinct first denominators.

For every odd divisor $d$ of $m_p$, the indicator of $d\equiv3\pmod4$ is
$(1-\chi_4(d))/2$.  Summation gives
\eqref{eq:diag-count-character}.  Multiplicativity gives
\[
 \sum_{d\mid m_p}\chi_4(d)
 =\prod_{q\equiv1(4)}(e_q+1)
  \prod_{r\equiv3(4)}\sum_{j=0}^{f_r}(-1)^j,
\]
and the last geometric sum is $\epsilon(f_r)$.  This proves
\eqref{eq:diag-count-factorized}.  The right side vanishes exactly when the
second prime product in \eqref{eq:diag-factorization} is empty, proving
\eqref{eq:diag-exceptional-exact}.
\end{proof}

\begin{corollary}[exact terminal multiplicity]
For $R\in\mathscr R_{\mathrm{diag}}(p)$ put
\begin{equation}\label{eq:diag-k}
 k_R=\frac{p+1}{R}\equiv2\pmod4.
\end{equation}
The corresponding terminal coordinates and matrix are
\begin{equation}\label{eq:diag-terminal-all}
 (\ell_0,\ell_1,\ell_2)
 =\left(\frac{k_R+2}{4},\frac{3k_R+2}{4},1\right),
 \qquad
 \mathsf R_R=
 \begin{pmatrix}
 3&-(3k_R+2)/4\\
 4&-(k_R+1)
 \end{pmatrix},
\end{equation}
with $\det\mathsf R_R=-1$.  These matrices are pairwise distinct as $R$
varies.  Hence $\nu_{\mathrm{diag}}(p)$ is simultaneously the number of
ordered diagonal-ray decompositions, diagonal-ray exterior selector states,
and determinant-$(-1)$ terminal matrices attached to this ray.
\end{corollary}

\begin{proof}
The formulas are \eqref{eq:diagonal-terminal} with $q$ replaced by the full
allowed divisor $R$.  Distinct $R$ give distinct $k_R=(p+1)/R$, and the
lower-right matrix entry is $-(k_R+1)$.
\end{proof}

\begin{corollary}[the exact hard residue split]
Every prime $p\equiv5\pmod8$ has at least one diagonal-ray certificate.  In
particular every hard prime $p\equiv13\pmod{24}$ is captured by this one ray.
A hard prime not captured by the ray must satisfy $p\equiv1\pmod{24}$ and
$(p+1)/2$ must lie in the multiplicative semigroup generated by primes
congruent to $1$ modulo four.
\end{corollary}

\begin{proof}
If $p\equiv5\pmod8$, then $m_p\equiv3\pmod4$, so its prime factorization
contains a prime congruent to $3$ modulo four.  Apply
\eqref{eq:diag-exceptional-exact}.  The final assertion follows by combining
that equivalence with the two possible hard classes modulo $24$.
\end{proof}


\begin{example}[the two diagonal certificates at $p=109$]
Here $m_p=55=5\cdot11$, and its divisors congruent to $3$ modulo four are
$11$ and $55$.  Formula \eqref{eq:diag-R-map} gives the two distinct ordered
solutions
\[
 \frac4{109}=\frac1{30}+\frac1{300}+\frac1{32700}
 =\frac1{41}+\frac1{82}+\frac1{8938}.
\]
Their terminal matrices are respectively
\[
 \begin{pmatrix}3&-8\\4&-11\end{pmatrix},
 \qquad
 \begin{pmatrix}3&-2\\4&-3\end{pmatrix},
\]
and both have determinant $-1$.
\end{example}

\section{Middle-channel rays}

For $n=p$ prime and $\delta=p$, \eqref{eq:projective-incidence} becomes
\begin{equation}\label{eq:middle-ray-condition}
 R\mid r+t.
\end{equation}
Thus a primitive ray admits a middle channel at precisely the divisors
\begin{equation}\label{eq:middle-R-set}
 \mathscr R_M(r,t)=\{R\mid r+t:R\equiv3\pmod4\}.
\end{equation}
This set is nonempty exactly when $r+t$ has a prime divisor congruent to
$3\pmod4$.

For fixed $R\in\mathscr R_M(r,t)$, put $A=4rt$, $g=\gcd(A,M)$ and
$M'=M/g$.

\begin{theorem}[middle-ray prime progressions]\label{thm:middle-progressions}
A prescribed prime class $p\equiv c\pmod M$ is compatible with this ray and
$R$ exactly when
\begin{equation}\label{eq:middle-compatibility}
 g\mid c+R.
\end{equation}
Under this condition the congruence
\begin{equation}\label{eq:middle-s-cong}
 \frac Ags_0\equiv\frac{c+R}{g}\pmod{M'}
\end{equation}
has one class modulo $M'$.  Put
\begin{equation}\label{eq:middle-p0Q}
 p_0=As_0-R,\qquad Q_M=AM'.
\end{equation}
Then $\gcd(p_0,Q_M)=1$, and every sufficiently large prime
$p\equiv p_0\pmod{Q_M}$ lies in the class $c\pmod M$ and gives
\begin{equation}\label{eq:middle-denominators}
 \boxed{
 x=rst,
 \qquad
 y=prs\,\frac{r+t}{R},
 \qquad
 z=pst\,\frac{r+t}{R},}
\end{equation}
where $s=(p+R)/(4rt)$.  These primes have asymptotic count
$\Li(X)/\varphi(Q_M)$.
\end{theorem}

\begin{proof}
The linear congruence $As\equiv c+R\pmod M$ is solvable exactly under
\eqref{eq:middle-compatibility}; division by $g$ gives
\eqref{eq:middle-s-cong}.  Since $R\mid r+t$ and $\gcd(r,t)=1$, one has
$\gcd(R,rt)=1$, so $\gcd(R,A)=1$.  The residue $p_0$ is a unit modulo $A$
because $p_0\equiv-R\pmod A$, and it is a unit modulo $M'$ because
$p_0\equiv c\pmod M$.  Therefore $\gcd(p_0,Q_M)=1$, and the prime number
theorem in arithmetic progressions applies.  For every such prime,
$p=4rts-R$ and $R\mid r+t$ give the middle incidence.  Substitution into the
source factor-coordinate reconstruction gives \eqref{eq:middle-denominators}.
\end{proof}

Unlike the exterior theorem, the allowed residuals are the finite set
\eqref{eq:middle-R-set} attached to the ray.

\section{Zariski density on the determinant cone}

Let
\begin{equation}\label{eq:det-cone}
 \mathscr C_{\det}=\{(u,a,v)\in\mathbb A^3:uv-a^2=0\}.
\end{equation}
Its positive rational projective rays are
\begin{equation}\label{eq:veronese}
 [r:t]\longmapsto[r^2:rt:t^2].
\end{equation}

\begin{theorem}[Zariski density of exterior certificate matrices]
\label{thm:zariski}
Fix $M,c$ as in \eqref{eq:M-c}.  The matrices
\begin{equation}\label{eq:certificate-matrices}
 \Gamma=s\mat{r^2}{rt}{rt}{t^2}
\end{equation}
arising from the exterior prime progressions of
\cref{thm:prime-saturation}, as $[r:t]$ ranges over primitive positive rays,
are Zariski dense in $\mathscr C_{\det}$.  Their images are Zariski dense in
the corresponding native and raw determinant-zero coefficient cones under
the invertible linear maps \eqref{eq:L-native} and \eqref{eq:raw-dict}.
\end{theorem}

\begin{proof}
Let $F(u,a,v)$ vanish on every constructed matrix.  Fix a primitive ray
$(r,t)$.  The prime-saturation theorem produces infinitely many distinct
positive contents $s$ on that ray, so
\[
 F(sr^2,srt,st^2)=0
\]
for infinitely many $s$.  This is a polynomial in $s$ and therefore vanishes
identically.  Hence every homogeneous component of $F$ vanishes at
$[r^2:rt:t^2]$.  The positive rational values of $t/r$ are Zariski dense in
$\Pone$, and the Veronese image \eqref{eq:veronese} is the irreducible
projective conic underlying \eqref{eq:det-cone}.  Thus every homogeneous
component vanishes on the cone, and so does $F$.  Invertible linear maps
preserve Zariski density.
\end{proof}

\section{Bézout openings of null rays}

The selector matrix \eqref{eq:Gamma-ray} is singular.  The next construction
adds an exact integral transverse direction instead of declaring the null ray
to be a regular period point.

Fix a primitive positive vector
\begin{equation}\label{eq:v-vector}
 \mathbf v=\begin{pmatrix}r\\t\end{pmatrix}.
\end{equation}
A Bézout companion is a vector
\begin{equation}\label{eq:w-vector}
 \mathbf w=\begin{pmatrix}x\\y\end{pmatrix},
 \qquad ry-tx=1.
\end{equation}
Then
\begin{equation}\label{eq:Bvw}
 B_{\mathbf v,\mathbf w}
 =\begin{pmatrix}r&x\\t&y\end{pmatrix}\in\SL_2(\Z).
\end{equation}

\begin{definition}[Bézout opening]\label{def:opening}
For $\delta>0$, define
\begin{equation}\label{eq:Aopening}
 \mathsf A_{\delta,\mathbf w}
 =\delta\mathbf v\mathbf v^{\mathsf T}
  -\mathbf w\mathbf w^{\mathsf T}
 =B_{\mathbf v,\mathbf w}
  \diag(\delta,-1)
  B_{\mathbf v,\mathbf w}^{\mathsf T}.
\end{equation}
\end{definition}

\begin{theorem}[integral opening and regular ambient period point]
\label{thm:opening}
Every opening satisfies
\begin{equation}\label{eq:opening-det}
 \det\mathsf A_{\delta,\mathbf w}=-\delta.
\end{equation}
The companions are exactly
$\mathbf w_k=\mathbf w+k\mathbf v$, $k\in\Z$, and
\begin{equation}\label{eq:opening-torsor}
 \mathsf A_{\delta,\mathbf w_k}
 =\mathsf A_{\delta,\mathbf w}
 -k(\mathbf v\mathbf w^{\mathsf T}
    +\mathbf w\mathbf v^{\mathsf T})
 -k^2\mathbf v\mathbf v^{\mathsf T}.
\end{equation}
The map $k\mapsto\mathsf A_{\delta,\mathbf w_k}$ is injective.

One may choose the companion so that $|x|\le r/2$.  For that choice,
\begin{equation}\label{eq:A11positive}
 (\mathsf A_{\delta,\mathbf w})_{11}
 =\delta r^2-x^2
 \ge\left(\delta-\frac14\right)r^2>0.
\end{equation}
Write
\begin{equation}\label{eq:Aentries}
 \mathsf A_{\delta,\mathbf w}=\mat{A_{11}}{A_{12}}{A_{12}}{A_{22}}.
\end{equation}
Then the raw coefficient triple
\begin{equation}\label{eq:raw-opening}
 (T_{\rm raw},M_{\rm raw},B_{\rm raw})
 =\left(A_{11},\frac{A_{12}}{\sqrt6},A_{22}\right)
\end{equation}
has
\begin{equation}\label{eq:raw-opening-disc}
 T_{\rm raw}B_{\rm raw}-6M_{\rm raw}^2=-\delta,
 \qquad
 T_{\rm raw}>0,
\end{equation}
and
\begin{equation}\label{eq:raw-sign}
 B_{\rm raw}-\frac{6M_{\rm raw}^2}{T_{\rm raw}}
 =-\frac\delta{A_{11}}<0.
\end{equation}
Thus the purely imaginary periods
\begin{equation}\label{eq:opening-periods}
 \tau=iA_{11},\qquad
 \mu=i\frac{A_{12}}{\sqrt6},\qquad
 \beta=iA_{22}
\end{equation}
give
\begin{equation}\label{eq:opening-ZP}
 Z_A=i\begin{pmatrix}
 \sqrt6A_{12}&A_{11}\\
 A_{22}&A_{12}/\sqrt6
 \end{pmatrix},\qquad
 P_A=\begin{pmatrix}
 0&0&1&0\\
 \sqrt6A_{12}&A_{11}&0&0\\
 0&0&0&1\\
 A_{22}&A_{12}/\sqrt6&0&0
 \end{pmatrix},
\end{equation}
with
\begin{equation}\label{eq:opening-P-det}
 \det P_A=A_{11}A_{22}-A_{12}^2=-\delta.
\end{equation}
Thus they give a full-rank lattice in the canonical ambient period form.  The native
coefficient triple is
\begin{equation}\label{eq:native-opening}
 (t,m,b)=\left(\frac{A_{11}}2,\frac{A_{12}}2,3A_{22}\right),
\end{equation}
for which
\begin{equation}\label{eq:native-opening-disc}
 tb-6m^2=-\frac32\delta,
 \qquad G(t,m,b)=\mathsf A_{\delta,\mathbf w}.
\end{equation}
The two triples are related exactly by \eqref{eq:D-map}.
\end{theorem}

\begin{proof}
Equation \eqref{eq:opening-det} follows from
$\det B_{\mathbf v,\mathbf w}=1$.  Every other solution of
$ry-tx=1$ differs by an integral multiple of $(r,t)$, and direct expansion
gives \eqref{eq:opening-torsor}.  In the basis $(\mathbf v,\mathbf w)$, the
opening has matrix
\[
 \begin{pmatrix}\delta-k^2&-k\\-k&-1\end{pmatrix},
\]
so equality for two integers $k$ forces equality of their off-diagonal
entries.  Shifting the companion chooses a centered representative of $x$
modulo $r$, proving \eqref{eq:A11positive}.

Equations \eqref{eq:raw-opening-disc} and \eqref{eq:raw-sign} are the
determinant identity divided by the positive entry $A_{11}$.  In the
canonical realification of
$Z=\bigl(\begin{smallmatrix}6\mu&\tau\\\beta&\mu\end{smallmatrix}\bigr)$,
the real period determinant is exactly the displayed determinant
\eqref{eq:opening-P-det}, equivalently
$T_{\rm raw}B_{\rm raw}-6M_{\rm raw}^2=-\delta\ne0$.  Therefore its four
columns span $\R^4$ and form a lattice.  Substitution into
\eqref{eq:native-G} gives \eqref{eq:native-opening-disc}; applying
\eqref{eq:D-map} gives \eqref{eq:raw-opening}.
\end{proof}

\begin{remark}[exact geometric scope]
The opening constructs a regular complex two-torus in the ambient canonical
period domain.  It supplies neither the real parts nor the holomorphic and
triangle-group equivariance conditions needed to place that point on the
global period image of the proposed $S^6$ family.  The obstruction is not
rank: the rank defect has been removed exactly.  The remaining question is
membership in that particular global period map.
\end{remark}

\begin{corollary}[opening is constant along a saturated prime ray]
Fix $(r,t)$, $\delta$, and a Bézout companion.  The matrix
$\mathsf A_{\delta,\mathbf w}$ is independent of the content $s$, the
external denominator $n$, and the residual $R$.  Thus every infinite prime
family produced on the fixed ray by \cref{thm:prime-saturation} maps to one
fixed regular ambient period point after the opening data are fixed.
\end{corollary}

\section{Exact left--right morphism to the terminal matrix}

Let $\mathsf R_\delta$ be the terminal matrix in
\eqref{eq:general-row-shear}.  Define
\begin{equation}\label{eq:P0Q0}
 P_0=\mat{-1}{1}{-4}{3},
 \qquad
 Q_0=\mat{1}{T_\delta-\ell_1}{0}{1}.
\end{equation}
Then
\begin{equation}\label{eq:terminal-diagonalization}
 P_0\mathsf R_\delta Q_0=\diag(1,-\delta).
\end{equation}
For $B=B_{\mathbf v,\mathbf w}$ and
$J=\mat{0}{1}{1}{0}$, put
\begin{equation}\label{eq:LA-RA}
 L_A=-JB^{-1},\qquad R_A=B^{-\mathsf T}J.
\end{equation}
Both are in $\GL_2(\Z)$ and
\begin{equation}\label{eq:opening-diagonalization}
 L_A\mathsf A_{\delta,\mathbf w}R_A=\diag(1,-\delta).
\end{equation}

\begin{theorem}[coordinate-level opening--terminal morphism]
\label{thm:GL2-morphism}
The terminal and opened matrices are related by the exact left--right
identity
\begin{equation}\label{eq:full-GL2-morphism}
 \boxed{
 \mathsf R_\delta
 =P_0^{-1}L_A\,
  \mathsf A_{\delta,\mathbf w}\,
  R_AQ_0^{-1}.}
\end{equation}
Every displayed left and right factor belongs to $\GL_2(\Z)$ and has an
explicit inverse.  Thus the relationship is stronger than equality of
determinants or Smith invariants.
\end{theorem}

\begin{proof}
A direct multiplication gives
\[
 P_0\mat{3}{-\ell_1}{4}{-T_\delta}
 =\mat{1}{\ell_1-T_\delta}{0}{4\ell_1-3T_\delta}
 =\mat{1}{\ell_1-T_\delta}{0}{-\delta}.
\]
Right multiplication by $Q_0$ gives
\eqref{eq:terminal-diagonalization}.  Equation \eqref{eq:Aopening} gives
\begin{align*}
 L_A\mathsf A_{\delta,\mathbf w}R_A
 &=-JB^{-1}B\diag(\delta,-1)B^{\mathsf T}B^{-\mathsf T}J\\
 &=-J\diag(\delta,-1)J=\diag(1,-\delta).
\end{align*}
Equating the two diagonalizations proves \eqref{eq:full-GL2-morphism}.
\end{proof}

\begin{example}[the fixed diagonal-ray opening]
For $\mathbf v=(1,1)^{\mathsf T}$ choose
$\mathbf w=(0,1)^{\mathsf T}$.  At $\delta=1$,
\begin{equation}\label{eq:Astar}
 \mathsf A_*=\mat{1}{1}{1}{0},\qquad \det\mathsf A_*=-1.
\end{equation}
The raw period triple is
\[
 (T_{\rm raw},M_{\rm raw},B_{\rm raw})
 =\left(1,\frac1{\sqrt6},0\right).
\]
For the density-one family \eqref{eq:diagonal-terminal}, every varying
terminal matrix is $\GL_2(\Z)$-equivalent through
\eqref{eq:full-GL2-morphism} to this one fixed ambient opening.
\end{example}

\section{Exact map to the defect-one Koide shell}

For an opening $A=\mathsf A_{\delta,\mathbf w}$, define its Lorentz coordinate
vector
\begin{equation}\label{eq:zA}
 \mathbf z(A)=
 \left(
 \frac{A_{11}+A_{22}}2,
 A_{12},
 \frac{A_{11}-A_{22}}2
 \right)=(H_A,X_A,Y_A).
\end{equation}
By construction,
\begin{equation}\label{eq:QAminusdelta}
 H_A^2-X_A^2-Y_A^2=\det A=-\delta.
\end{equation}
Define the exact homothety and orthogonal inverse
\begin{equation}\label{eq:defect-map}
 \mathfrak K_\delta(A)
 =\mathcal F^{-1}\left(\frac{\mathbf z(A)}{\sqrt{3\delta}}\right).
\end{equation}

\begin{theorem}[channel opening to the defect-one Koide hyperboloid]
\label{thm:defect-one-map}
For every divisor channel order $\delta>0$ and every Bézout opening,
\begin{equation}\label{eq:QKminus-third}
 Q_K(\mathfrak K_\delta(A))=-\frac13.
\end{equation}
Equivalently, with
\begin{equation}\label{eq:Delta}
 \Delta(q)=4(q_1q_2+q_1q_3+q_2q_3)-(q_1^2+q_2^2+q_3^2)=3Q_K(q),
\end{equation}
one has
\begin{equation}\label{eq:Deltaminus1}
 \boxed{\Delta(\mathfrak K_\delta(A))=-1.}
\end{equation}
The map is defined over
$\Q(\sqrt2,\sqrt3,\sqrt\delta)$ for integral $A$.  It need not land in the
integral defect-one shell.
\end{theorem}

\begin{proof}
The map $\mathcal F$ is orthogonal and identifies $Q_K$ with the Lorentz form.
Equations \eqref{eq:QAminusdelta} and \eqref{eq:defect-map} therefore give
\[
 Q_K(\mathfrak K_\delta(A))
 =\frac{-\delta}{3\delta}=-\frac13.
\]
The identity $\Delta=3Q_K$ gives \eqref{eq:Deltaminus1}.  The displayed
radicals in $\mathcal F^{-1}$ and in the homothety give the stated field.
\end{proof}


\begin{theorem}[exact rational nonintersection of the orthogonal defect map]
\label{thm:defect-map-no-rational-points}
Let $\delta\in\mathbb Q_{>0}$ and let
$A\in\operatorname{Sym}_2(\mathbb Q)$ satisfy $\det A=-\delta$.  Then
\begin{equation}\label{eq:no-rational-image}
 \boxed{\mathfrak K_\delta(A)\notin\mathbb Q^3.}
\end{equation}
Consequently the image of the integral Bézout openings under the exact map
\eqref{eq:defect-map} is disjoint from the integral defect-one shell, not
merely unproved to lie in it.
\end{theorem}

\begin{proof}
Write $\mathbf z(A)=(H_A,X_A,Y_A)$ and suppose
$q=\mathfrak K_\delta(A)=(q_1,q_2,q_3)$ were rational.  Put
\begin{equation}\label{eq:SUV}
 S=q_1+q_2+q_3,
 \qquad U=q_1-q_3,
 \qquad V=q_1+q_3-2q_2.
\end{equation}
Applying the exact orthogonal map \eqref{eq:F} to
$q=\mathcal F^{-1}(\mathbf z(A)/\sqrt{3\delta})$ gives
\begin{equation}\label{eq:rationality-relations}
 H_A=\sqrt\delta\,S,
 \qquad
 X_A=\sqrt{\frac{3\delta}{2}}\,U,
 \qquad
 Y_A=\sqrt{\frac\delta2}\,V.
\end{equation}
All six quantities $H_A,X_A,Y_A,S,U,V$ are rational.

If two among $H_A,X_A,Y_A$ are nonzero, division of the corresponding two
relations in \eqref{eq:rationality-relations} would make one of
$\sqrt{3/2}$, $1/\sqrt2$, or $\sqrt3$ rational.  Hence at most one of
$H_A,X_A,Y_A$ is nonzero.  The case in which only $H_A$ is nonzero is
impossible because
$H_A^2-X_A^2-Y_A^2=-\delta<0$.  If only $X_A$ is nonzero, then
$X_A^2=\delta$ and the second relation gives $U^2=2/3$, impossible for
$U\in\mathbb Q$.  If only $Y_A$ is nonzero, then $Y_A^2=\delta$ and the
third relation gives $V^2=2$, again impossible over $\mathbb Q$.  The zero
triple contradicts $\det A=-\delta$.  Thus no rational $q$ exists.
\end{proof}

\begin{remark}[exact separation from the counted integral orbit]
The map \eqref{eq:defect-map} proves equality of the quadratic level
$\Delta=-1$ and produces a real algebraic point for every opening.  Theorem
\ref{thm:defect-map-no-rational-points} proves that this particular map has
empty intersection with $\mathbb Q^3$, hence with the integral chamber counted
in the Koide defect-one theorem.  Any arithmetic transfer to that counted
orbit requires a different, explicitly constructed morphism.
\end{remark}

\begin{remark}[why the existing count does not transfer]
The current Koide arithmetic workbench counts integral points on a particular
reduced chamber of $\Delta=-1$ and proves a linear asymptotic there.  The map
\eqref{eq:defect-map} produces exact algebraic points and, by
\cref{thm:defect-map-no-rational-points}, none of them is rational.  Therefore
the proved integral counting constant is not a count of Erdős--Straus selector
states.  What is transferred is the
quadratic level and the complete coordinate map, not integrality or chamber
membership.
\end{remark}

\section{Faithful certificates and exact inverse data}

An oriented divisor channel determines the pair
\begin{equation}\label{eq:double-cert}
 \mathfrak C(n,R,\delta,u,v)
 =\left(
 \Gamma=\mat{u}{a}{a}{v},
 \quad
 \mathsf R_\delta=\mat{3}{-\ell_1}{4}{-T_\delta}
 \right).
\end{equation}
The two determinant equations are
\begin{equation}\label{eq:two-dets}
 \det\Gamma=0,
 \qquad
 \det\mathsf R_\delta=-\delta,
\end{equation}
and the coupling equations are
\begin{equation}\label{eq:coupling}
 4a=n+R,
 \qquad
 RT_\delta=4u+\delta,
 \qquad
 uv=a^2.
\end{equation}

\begin{theorem}[faithfulness of the double certificate]
\label{thm:faithful}
The map from marked oriented channel states to the pair
\eqref{eq:double-cert} is injective.  The inverse is
\begin{align}
 u&=\Gamma_{11},&a&=\Gamma_{12},&v&=\Gamma_{22},\label{eq:inverse-Gamma}\\
 T_\delta&=-(\mathsf R_\delta)_{22},&
 \ell_1&=-(\mathsf R_\delta)_{12},\label{eq:inverse-R}\\
 j&=\frac{u-\ell_1}{T_\delta},&
 R&=4j+3,\label{eq:inverse-j}\\
 n&=4a-R,&
 \delta&=-\det\mathsf R_\delta.\label{eq:inverse-n-delta}
\end{align}
The terminal matrix alone is not faithful; the rank-one matrix restores the
external denominator, shell marker, divisor support, and channel orientation.
\end{theorem}

\begin{proof}
Equations \eqref{eq:inverse-Gamma} and \eqref{eq:inverse-R} read off matrix
entries.  The definition $\ell_1=u-jT_\delta$ gives
\eqref{eq:inverse-j}; then \eqref{eq:aN} and
\eqref{eq:terminal-snf} give \eqref{eq:inverse-n-delta}.  Substitution
reconstructs every original coordinate.
\end{proof}

\section{Exact consequences and remaining boundary}

The proved chain is
\begin{equation}\label{eq:full-chain}
 \boxed{
 \begin{gathered}
 \text{oriented fixed-shell reciprocal state}
 \longleftrightarrow
 (\delta,u,v),\\
 (\delta,u,v)
 \longmapsto
 \left(\Gamma,\mathsf R_\delta\right),\\
 \Gamma\longmapsto[\mathbf v],
 \qquad
 ([\mathbf v],\delta,\mathbf w)
 \longmapsto\mathsf A_{\delta,\mathbf w},\\
 \mathsf A_{\delta,\mathbf w}
 \xrightarrow{\ \GL_2(\Z)\times\GL_2(\Z)\ }
 \mathsf R_\delta,\\
 \mathsf A_{\delta,\mathbf w}
 \longmapsto
 \text{regular ambient period point},\\
 \mathsf A_{\delta,\mathbf w}
 \xrightarrow{\ \mathfrak K_\delta\ }
 \{\Delta=-1\}.
 \end{gathered}}
\end{equation}
Every arrow has been displayed in coordinates and supplied with its inverse
when it is bijective.

The following statements are exact consequences.
\begin{enumerate}[label=\textup{(\roman*)}]
\item The prime two-channel selector is the specialization of a complete
$\delta\mid n^2$ channel theory.
\item Complementary denominator orientations correspond to reciprocal
terminal orders whose product is $n^2$, and to opposite Lorentz rapidities.
\item Every primitive positive determinant-null ray supports infinitely many
exterior prime certificates in every admissible prime progression.
\item One ray, $[1:1]$, supports relative density one of those primes; its
certificate multiplicity for each individual prime is exactly the divisor
count in \eqref{eq:diag-count-character}, and it carries a fixed
determinant-$(-1)$ ambient opening.
\item Every channel ray has infinitely many Bézout openings of determinant
$-\delta$, and each opening is connected to its terminal matrix by the exact
left--right identity \eqref{eq:full-GL2-morphism}.
\item Every opening maps to the real defect-one Koide hyperboloid, but the exact
orthogonal--homothetic image is disjoint from $\mathbb Q^3$; no integral point
count is transferred by this morphism.
\end{enumerate}

The Erdős--Straus conjecture remains the surjectivity statement in the other
quantifier direction: every prime must lie in at least one selector fibre.
Neither density one, Zariski density, nor infinite saturation of every fixed
ray rules out a density-zero set of primes missed by all rays.

\section*{Computational receipt}

The accompanying exact-arithmetic verifier checks:
\begin{itemize}
\item the channel bijection and fixed-point-free complement in $665$ coprime
shells with $n\le161$;
\item $1430$ oriented channel states and $512$ explicit
opening--terminal $\GL_2(\Z)$ morphisms;
\item the prime specialization in $9137$ shells for primes $p<1000$;
\item $57$ character-sum channel counts;
\item $192$ primitive-ray/congruence-class progressions and explicit prime
certificates in all $48$ reduced classes $c\equiv1\pmod{12}$ modulo $840$ on
the diagonal ray;
\item $30$ Bézout-opening torsors with eleven consecutive companion shifts
each;
\item the complete diagonal-ray divisor count for every prime $p<5000$ with
$p\equiv1\pmod4$, including the exact terminal matrix multiplicity.
\end{itemize}
These checks are bounded receipts for the displayed identities.  The proofs
above establish the unbounded statements.

\begin{thebibliography}{99}
\bibitem{AlpogeFable2026}
L. Alpöge, with disclosed Claude/Fable contribution,
\emph{The $(3,4,\infty)$ modular family of $2$-tori, completed at its three
special points, is a complex structure on $S^6$}, 2026; also circulated as
\emph{A compact complex threefold fibred by tori over the projective line,
and the six-sphere}.  Citable research record:
\url{https://doi.org/10.5281/zenodo.22238779}.

\bibitem{ClankersES2026}
The Clankers,
\emph{Erdős--Straus Project Archive: Residual Divisor Shells, Quotient
Surpluses, and Finite Verification}, version 2026.08.31.1, Zenodo (2026),
\url{https://doi.org/10.5281/zenodo.22209908}.

\bibitem{ErdosStraus1948}
P. Erdős and E. G. Straus,
problem on the decomposition of $4/n$ into three unit fractions, 1948.

\bibitem{Koide1983}
Y. Koide,
\emph{New view of quark and lepton mass hierarchy},
Physical Review D \textbf{28} (1983), 252--254,
\url{https://doi.org/10.1103/PhysRevD.28.252}.

\bibitem{Kocik2012}
J. Kocik,
\emph{The Koide lepton mass formula and geometry of circles},
arXiv:1201.2067 (2012).

\bibitem{Davenport}
H. Davenport,
\emph{Multiplicative Number Theory}, 3rd ed., revised by H. L. Montgomery,
Springer, 2000.
\end{thebibliography}

\end{document}
